\documentclass[11pt]{article}
\usepackage[margin=25mm]{geometry}
\usepackage{fontspec}
\setmainfont{Latin Modern Roman}
\usepackage{amsmath,amssymb,mathtools}
\usepackage{unicode-math}
\setmathfont{Latin Modern Math}
\usepackage{xurl,hyperref}
\hypersetup{colorlinks=true,urlcolor=blue}
\setcounter{secnumdepth}{0}
\setlength{\parindent}{0pt}
\setlength{\parskip}{5pt}
\setlength{\emergencystretch}{3em}
\newcommand{\sref}[2]{\hyperlink{source:#1:#2}{[S#2]}}
\newcommand{\cataloguescope}[1]{\paragraph{Recorded target.}#1}
\newcommand{\R}{\mathbb R}
\newcommand{\C}{\mathbb C}
\newcommand{\N}{\mathbb N}
\newcommand{\Z}{\mathbb Z}
\newcommand{\Q}{\mathbb Q}
\newcommand{\F}{\mathbb F}
\newcommand{\eps}{\varepsilon}
\begin{document}
% UnsolvedMath ID 7900003; source catalogue code AMR-078-0003
\setcounter{section}{7900002}
\section{Quantum Diffusion in the Anderson Model}\label{7900003}
{\small UnsolvedMath ID 7900003 · Mathematical Physics}
\subsection{Definitions and mathematical statement}

\paragraph{Shared notation.}$\mathbb N=\{1,2,\ldots\}$; the standard number fields have their usual meanings. All additional conventions are stated locally.


\paragraph{Definitions and notation.}On $\ell^2(\mathbb Z^\nu)$, the source defines $(Hu)(n)=\sum_{|j|=1}u(n+j)+V_\omega(n)u(n)$, where the site potentials are i.i.d. uniform on $[a,b]$. The matrix element $e^{itH}(n,0)$ is $\langle\delta_n,e^{itH}\delta_0\rangle$, and $n^2$ means $|n|^2$. Thus $M_2(t)=\sum_n|n|^2|e^{itH}(n,0)|^2$ is the second moment of position for a particle initially at the origin. The target is diffusive $M_2(t)\sim ct$, conditional on an absolutely continuous regime.

\paragraph{Catalogue wording retained for traceability.}
For the Anderson model in dimension $\nu\geq3$ and disorder strengths $|b-a|$ admitting absolutely continuous spectrum, prove that $\sum_{n\in\mathbb Z^\nu} n^2|e^{itH}(n,0)|^2$ grows asymptotically like $ct$ as $t\to\infty$.
\paragraph{Scope and source qualification.}The statement is conditional on a regime with absolutely continuous spectrum, whose existence is a related problem. The catalogue does not specify whether the asymptotic is almost sure, averaged over disorder, or energy-restricted; that ambiguity is retained rather than replaced by one stronger assertion.
\subsection{Short English statement}
Show that a particle in the conducting regime of the Anderson model spreads diffusively.
\subsection{Research attempt}
\paragraph{Model, date and investigated scope.}
OpenAI GPT-6 Astra (configured model \texttt{gpt-6-astra}), ultra reasoning; generated 11 October 2026. Computed exact ballistic spreading at zero disorder and a ballistic upper bound at bounded disorder; no positive-disorder diffusion theorem. The source ambiguity in averaging and the zero-disorder endpoint remain explicit.

\paragraph{Sources and literature check.}
Read the supplied complete definition, catalogue wording, scope qualifications and source list for this canonical ID. The original-source citations are inherited as supplied references; no fresh primary-source or novelty search was performed for this attempt. No current open-status claim is inferred from this calculation.

\paragraph{The zero-disorder obstruction.}Take the endpoint case \(a=b=v_0\), a deterministic potential. On the Fourier torus \([-\pi,\pi]^\nu\), the Hamiltonian is multiplication by \(h(k)=v_0+2\sum_{j=1}^\nu\cos k_j\), and the initial vector \(\delta_0\) transforms to the constant function one with normalized Haar measure. Position multiplication corresponds to differentiation in \(k\). Therefore Parseval's identity gives the exact formula
\[M_2(t)=\sum_j\int\left|\partial_{k_j}e^{ith(k)}\right|^2\frac{dk}{(2\pi)^\nu}=4t^2\sum_j\int\sin^2k_j\frac{dk}{(2\pi)^\nu}=2\nu t^2.\]
The spectral measures in this model are absolutely continuous: a single coordinate variable \(2\cos k_1\) has density \(1/(\pi\sqrt{4-E^2})\) on \((-2,2)\), and adding the other independent cosine variables convolves this density with probability measures. Weighted Fourier spectral measures are also absolutely continuous, by absolute continuity of their measure with respect to Haar measure and its energy pushforward. Thus absolutely continuous spectrum alone does not force diffusion.
\paragraph{A robust upper bound for nonzero disorder.}For any bounded diagonal potential, the commutator \([X_j,H]\) is a difference of two unit shifts and has norm at most two. Integrating the Heisenberg equation on \(\delta_0\), whose position is zero, gives \(\|X_je^{itH}\delta_0\|\le2|t|\), and hence \(M_2(t)\le4\nu t^2\). Finite-range hopping justifies this on truncated position operators followed by the usual monotone moment limit; alternatively the power series of the bounded Hamiltonian provides finite moments at each fixed time and the same integrated commutator identity.
\paragraph{Remaining target.}If the intended disorder is strictly positive, the zero-disorder calculation is a boundary countercheck, not a disproof of that intended assertion. If zero is literally included, it excludes an unrestricted conditional formulation. Neither calculation gives the linear leading term for positive random disorder. The first missing step is a quantitative long-time decorrelation estimate for the velocity, with the required disorder-averaged, almost-sure or energy-restricted meaning fixed. Those meanings are not chosen silently here.

\paragraph{Verification and adversarial audit.}
Checked Fourier normalization and coefficient 2 nu, absolute continuity of the free energy pushforward, commutator norms and the zero-disorder scope distinction. This was an LLM adversarial audit, not independent human verification; no theorem prover was used. Prompts A then B applied to this ID: Checked Fourier normalization and coefficient 2 nu, absolute continuity of the free energy pushforward, commutator norms and the zero-disorder scope distinction. Scope remains limited: Prove velocity decorrelation sufficient for a linear second-moment asymptotic at strictly positive disorder, in a specified probabilistic/energy sense.

Coordinator checked the Fourier second-moment identity, absolute continuity of the free spectral measures and bounded hopping-commutator estimate. The zero-disorder endpoint is distinguished from positive disorder. This additional check was model-assisted, not independent human verification.

\paragraph{Tools, process and contributions.}
Python 3 and local shell used to read supplied files and serialize this manuscript; mathematics derived symbolically in the text, with no numerical experiment or formal proof assistant. Contribution prompt A was followed by prompt B. GPT-6 Astra with ultra reasoning drafted the research and performed its stated self-audit; any additional coordinator contribution is identified in the verification and audit. A GPT-6-based Codex coordinator assembled the manuscript and checked repository integration; its exact variant and reasoning configuration were not exposed. No human mathematical contribution, independent human proof check or formal proof certification is supplied. Complete derivations and the consulted source references are included here; private reasoning transcripts are not published.

\paragraph{Final status and first remaining gap.}
\textbf{UNRESOLVED}. Prove velocity decorrelation sufficient for a linear second-moment asymptotic at strictly positive disorder, in a specified probabilistic/energy sense.

\section{Completion Estimate}
\noindent\textbf{COMPLETION: 0\%}
The calculation is an elementary benchmark or conditional reduction; it establishes no claimed new portion of the complete source target.

\subsection{Sources}
\begin{itemize}
\item[\hypertarget{source:7900003:1}{S1}] ULAM.AI and the UnsolvedMath Contributors, UnsolvedMath, record 7900003 (AMR-078-0003). Catalogue curation under CC BY 4.0; definitions and English exposition added locally.
\url{https://huggingface.co/datasets/ulamai/UnsolvedMath}.
\item[\hypertarget{source:7900003:2}{S2}] Barry Simon, Schrödinger operators in the twenty-first century, Mathematical Physics 2000, pp. 283–288, original Problem 3.
\url{https://web.ma.utexas.edu/mp_arc-bin/mpa?yn=00-78}.
\item[\hypertarget{source:7900003:3}{S3}] Wikipedia contributors, Simon problems, pinned revision 1335881350, The 2000 list. Existing catalogue wording adapted; CC BY-SA 4.0.
\url{https://en.wikipedia.org/w/index.php?title=Simon_problems&oldid=1335881350#The_2000_list}.
\end{itemize}
\label{end:7900003}
\end{document}
